% Sources: author

\section{Discussion}
\label{sec:discussion}

% OUTLINE: Physics left out: Comptonization; relativistic effects (Newtonian disc, eta_acc = 1/(2 r_in_factor)).
% OUTLINE: Focusing by the funnel walls and the oblique view through the wind: corrected in the model
%          (Sect. sec:spectrum:transfer, mc_transfer = 1); checks and limits in Appendix app:mc
%          (largest errors beyond theta_c: BH up to -22/+23%, NS dipole model faint by up to x1.6).
% OUTLINE: Limits of the fixed choices (velocity law, R_sph and R_trap rules).
% OUTLINE: Energy carried by the wind through the photosphere (kinetic + enthalpy flux,
%          old manuscript "Wind power") vs eps_w (L_disc - L_disc,seen); diagnostic
%          wind_energy_loss_through_green_surface() in psrcs/wind_emission.py (run_many.py),
%          not part of the fitted model.

\subsection{Does the wind emit a blackbody?}
\label{sec:disc:thermalization}
% Source: ULXspectra-mc docs/thermalization_in_wind.md (PLAN C21), branch mc-transfer-corrections,
%         commit 69b9de7; estimate script docs/scripts/mc_table/thermalization.py

The model gives each disc ring a blackbody and each ring of the wind photosphere ($\taup = 1$) a diluted blackbody at the colour temperature of the disc that feeds it (Sect.~\ref{sec:spectrum}).
The Monte Carlo transfer of Appendix~\ref{app:mc} scatters grey and coherently, without absorption or emission, so each emitter keeps its spectral shape as well.
A blackbody, or a modified blackbody, at the temperature of the wind would require that the photons be made by absorption in the wind itself.
Here we estimate whether the wind can do that; the result is the reason for the dilution.
% source: thermalization_in_wind.md, sect. 1

At each interaction a photon is absorbed with the probability
\begin{equation}
\varepsilon_\nu = \frac{\alpha_\nu^{\rm abs}}{\alpha_\nu^{\rm abs} + n_{\rm e}\sigma_{\rm T}} .
\label{eq:eps_abs}
\end{equation}
% source: thermalization_in_wind.md, sect. 2; code: docs/scripts/mc_table/thermalization.py (eps)
A photon that starts at the scattering optical depth $\tau$ random-walks out in about $\tau^2$ scatterings, and its chance of being absorbed on the way is about $\varepsilon_\nu\tau^2$.
Absorption, and so thermal emission, is therefore likely deeper than the thermalisation depth
\begin{equation}
\tau_*(\nu) \approx \frac{1}{\sqrt{3\varepsilon_\nu}} ,
\label{eq:tau_therm}
\end{equation}
% source: thermalization_in_wind.md, sect. 2; Zeldovich & Shakura (1969), Rybicki & Lightman (1979) as named there
measured in scattering optical depth \citep{ZeldovichShakura1969, Shakura1973, RybickiLightman1979}.
This is the depth that the medium must have.
If the medium is deeper than $\tau_*$, photons are made near $\tau_*$ and scatter out, and the result is a modified blackbody at the temperature of that layer.
If the whole column is shallower, $\tau_{\rm col} < \tau_*$, the medium is effectively optically thin: thick to scattering but thin to absorption.
It absorbs a fraction $\varepsilon_\nu\tau_{\rm col}^2 \approx (\tau_{\rm col}/\tau_*)^2/3$ of the photons crossing it, and by Kirchhoff's law its own thermal emission is the same small fraction of a blackbody; the light that leaves it was made elsewhere and scattered.
For example, $\taup = 3$ and $\tau_* = 100$ give $\varepsilon_\nu \approx 3\times10^{-5}$; a crossing takes about 9 scatterings, so a photon is absorbed with the probability of only about $3\times10^{-4}$.
% source: thermalization_in_wind.md, sect. 2 (example)

To estimate $\tau_*$ in the model's wind, we take free-free absorption only, in fully ionised hydrogen with the Gaunt factor of unity and the gas temperature equal to the temperature of the wind ring:
\begin{equation}
\alpha_\nu^{\rm ff} = 3.7\times10^{8}\,T^{-1/2}\,n_{\rm e}n_{\rm i}\,\nu^{-3}\left(1 - {\rm e}^{-h\nu/kT}\right)\ {\rm cm^{-1}}
\label{eq:alpha_ff}
\end{equation}
% source: thermalization_in_wind.md, sect. 3; code: thermalization.py (a = 3.7e8 T^-0.5 ne^2 nu^-3 (1 - exp(-h nu/kT))), cgs
(cgs units).
We evaluate $\tau_*$ at the density of the photosphere and compare it with the whole vertical column of the wind, $\tau_{\perp,{\rm tot}}$, in BH3a fits of \srcHoIX and \srcNGCfiftyfour with $\omf = 0.7$ (Table~\ref{tab:therm}).
In \srcNGCthirteen at $\omf = 0.7$ no wind rings radiate.
% source: thermalization_in_wind.md, sect. 3 ("BH3a calibration fits at omega = 0.7": the BH3a full fits with omega_frac
%         = 0.7, docs/scripts/mc_table/runs.txt); commit 69b9de7

\begin{table}
\centering
\caption{Thermalisation depth $\tau_*$ (Eq.~\ref{eq:tau_therm}, free-free absorption), luminosity-weighted median over the wind rings, and the share of the wind light of \srcNGCfiftyfour that is thermalised. The whole vertical column is $\tau_{\perp,{\rm tot}} = 1.1$--$2.8$ in \srcHoIX and 1.4--33 in \srcNGCfiftyfour.}
\label{tab:therm}
\small
\setlength{\tabcolsep}{3pt}
\begin{tabular}{@{}cccc@{}}
\toprule
& \multicolumn{2}{c}{$\tau_*$} & thermalised \\
$E$ [keV] & Ho IX X-1 & NGC 5408 X-1 & NGC 5408 X-1 \\
\midrule
0.01 & 3.1 & 2.0 & 0.95 \\
0.1 & 47 & 23 & 0.23 \\
0.3 & 230 & 89 & 0.01 \\
1 & 1400 & 460 & 0 \\
3 & 7200 & 2300 & 0 \\
\bottomrule
\end{tabular}
% source: thermalization_in_wind.md, sect. 3 table; ULXspectra-mc commit 69b9de7
\end{table}

Above about 0.1--0.3~keV, $\tau_*$ far exceeds the column, so the wind cannot make X-ray photons.
Its X-rays are disc photons, made where the disc is dense enough to thermalise them, and scattered by the wind.
Scattering conserves the photon energy, so they keep the colour temperature of the disc under the wind; spread over the larger photosphere, they form a diluted blackbody, not a blackbody at the effective temperature of the wind.
In the UV, near 0.01~keV, $\tau_* \approx 2$--$3$ is below the column of \srcNGCfiftyfour, the wind thermalises, and a modified blackbody at the gas temperature applies.
% source: thermalization_in_wind.md, sect. 3

These estimates are rough.
Deeper in the column the density is higher, and $\tau_*$ is about 1.5 times lower (30 instead of 47 at 0.1~keV for \srcHoIX).
Bound-free absorption by metals raises the opacity at 0.1--1~keV and lowers $\tau_*$ by a factor of 3--10, which does not change the conclusion above about 0.3~keV.
Compton energy exchange is weak: $y = (4kT_{\rm e}/m_{\rm e}c^2)\,N \approx 0.07$ for $T_{\rm e} = 10^6$~K and $N = 100$ scatterings.
Radiation trapped in the outflow ($\tau v/c > 1$) cools adiabatically, which shifts photon energies without thermalising them.
Finally, the gas temperature of the wind is not known in the model.
% source: thermalization_in_wind.md, sect. 3, caveats

The same reasoning applies to the spherical envelope of \citet{Lipunova1999}, whose optical spectrum is a modified blackbody \citep{ZeldovichShakura1969}.
For the parameters of SS~433 used there, optical photons thermalise at $\tau_* \approx 2$ inside the envelope (Appendix~\ref{app:envelope}), so the modified blackbody is correct there.
From the optical at $3.5\times10^4$~K to 1~keV at $10^6$~K, however, the free-free coefficient decreases by a factor of about $2.5\times10^8$, and $\tau_* \propto (\alpha^{\rm ff}_\nu)^{-1/2}$ grows by about $1.6\times10^4$.
Moreover, the wind of our model is a finite column, $\taup \approx 1$--$30$ above the disc, not an envelope with $n \propto r^{-2}$ reaching ever denser layers inwards.
% source: thermalization_in_wind.md, sect. 4 (last paragraph): k_ff falls by ~4e-9, tau_* grows by ~1.6e4

A physically consistent spectrum of the wind therefore changes regime with frequency.
Where $\tau_*(\nu)$ is below the column of the wind, the wind thermalises and emits a modified blackbody at the gas temperature: the UV of ULX winds, as the optical of the envelope of SS~433.
Where $\tau_*(\nu)$ exceeds the column, the emission of the wind is scattered disc light that keeps the colour temperature of the disc: the X-rays above about 0.3~keV.
The model takes the second regime for the whole spectrum of the wind (Eq.~\ref{eq:Inu_wind}), since the fits are made to the X-rays.
% source: README "Wind spectrum"; PLAN C21 (author's decision 2026-11-04: option (b), diluted wind, wind_spectrum = 1)
In the UV this is not correct: there a blackbody at $\Tw$ would be the better description.
A spectrum that switches regime with frequency needs the gas temperature of the wind, which the model does not have.
The slope of the SEDs below 0.3~keV (Sect.~\ref{sec:fitting:data}) is an extension of a phenomenological fit and does not constrain this.
% source: thermalization_in_wind.md, sect. 5; README "Wind spectrum" (limitations)
% AUTHOR: the disc rings are also pure blackbodies (no colour correction; PLAN C21 quotes f_col ~ 1.6-2
%         for X-ray discs) - discuss here as well?
